\documentclass[10pt]{article}
\usepackage[margin=1.6cm]{geometry}
\usepackage{graphicx}
\usepackage{amsmath}
\usepackage{amssymb}
\usepackage[hidelinks]{hyperref}
\setlength{\parskip}{2pt}\setlength{\parindent}{0pt}

\title{\vspace{-1.4cm}Self-regulating accessibility: the unfrozen $d2{+}d4$ network never
enters the compositional phase\\
\large interpolation from frozen-$d2$ to fully simultaneous training --- 7 pretraining
durations $\times$ 8 global $\lambda$ $\times$ 3 seeds}
\author{}\date{\vspace{-1.2cm}}

\begin{document}\maketitle\vspace{-1.9cm}

\section*{Design}
Same 4-bit, full-cube, two-block ReLU residual setup as the frozen experiments
($N{=}1024$, init variance $1/N$, exact Walsh, full-batch SGD). \textbf{Phase 1} pretrains
\emph{everything} (embedding, both blocks, 3-output readout with the $d4$ head masked) on the
parents $x_1x_2,\,x_3x_4$ with $\lambda=0$, saving checkpoints (parameters \emph{and} momentum
buffers). \textbf{Phase 2} resumes from each checkpoint with $d4$ supervision on, the $d2$
losses kept, \emph{all layers trainable}, and a global weight decay
$\lambda\in\{0,10^{-4},3{\cdot}10^{-4},10^{-3},3{\cdot}10^{-3},10^{-2},3{\cdot}10^{-2},10^{-1}\}$
on all four weight matrices; within a seed every $\lambda$ starts from the identical training
state. Circuits are classified from block-2 pre-activation Walsh shares as before.

A design fact discovered at launch: \textbf{$R_0=Q_2/Q_1$ at the input to block 2 is not freely
dialable when nothing is frozen}. Random ReLU features already give $R_0\approx0.19$ at
initialization, and converged $d2$ pretraining saturates at $R_0\approx0.33$ --- the
$1{+}1\!\to\!2$ corner units emit singletons of the same per-neuron magnitude as their pair, so
$Q_1$ grows in lockstep with $Q_2$, capping $R_0$ \emph{below} $R_c=1/\sqrt8\approx0.354$
(with a brief overshoot to $0.34$ at $t_{\rm pre}\!\approx\!40$). Pretraining \emph{duration}
($t_{\rm pre}\in\{0,5,15,40,100,400,10^4\}$, from fully simultaneous to fully converged) is
therefore the controlled axis, and $R_0$ the measured covariate.

\section*{Findings}
\textbf{1. The direct circuit wins everywhere.} Across all 168 cells the compositional fraction
never exceeds $f_{2+2}=0.084$, while $d4$ is always solved ($\min|\hat f_4|=0.85$, no
failures). The $(R_0,\lambda)$ phase diagram (Fig.~\ref{fig:phase}, left) is uniformly in the
$1^4$ phase.

\textbf{2. Pretraining history washes out.} At $\lambda=0$, $f_{2+2}$ moves only
$0.044\to0.060$ between $t_{\rm pre}=0$ and $t_{\rm pre}=10^4$, although $R_0$ spans
$0.19\to0.34$. All conditions converge to the \emph{same} $R_{\rm final}\approx0.31$: whatever
the starting accessibility, joint training drives the representation to one fixed point below
$R_c$ (Fig.~\ref{fig:plane} --- trajectories from both directions meet in the same band,
never approaching $R_c$).

\begin{figure}[h!]\centering
\includegraphics[width=0.495\textwidth]{phase_R0.pdf}\hfill
\includegraphics[width=0.495\textwidth]{phase_Rfinal.pdf}\vspace{-4pt}
\caption{Phase diagrams of the unfrozen experiment vs.\ $R_0$ (left) and $R_{\rm final}$
(right); color = mean $f_{2+2}$ on the same scale as the frozen-case diagram. Everything stays
in the $1^4$ (blue) phase, left of $R_c$ (red). Note in the right panel how increasing
$\lambda$ \emph{lowers} $R_{\rm final}$ while all pretraining durations collapse onto nearly
the same point per $\lambda$.}
\label{fig:phase}
\end{figure}

\textbf{3. Global decay lowers accessibility --- the opposite of its frozen-case role.} With
$\lambda$ acting on $V^{(1)}$ too, decay prunes the corner machinery to the minimum that still
solves $d2$: $\|V^{(1)}\|_F$ falls $32\to1.8$ and $R_{\rm final}$ falls $0.32\to0.21$ as
$\lambda\!:0\to10^{-1}$ --- an inverse $\alpha$-rescaling. Consequently $f_{2+2}$ hits exactly
$0$ by $\lambda=3{\cdot}10^{-3}$. In the frozen setup $\lambda$ sharpened the compositional
phase right of $R_c$; here it destroys the conditions for it.

\textbf{4. Frozen and unfrozen collapse against $R_{\rm final}$, not against $R_0$.} At matched
$(R_{\rm final},\lambda)$ the two experiments agree quantitatively --- e.g.\ at $\lambda=0$,
unfrozen $f_{2+2}=0.060$ at $R_{\rm final}=0.32$ vs.\ frozen $0.052$ at $R=0.32$; at
$\lambda\geq3{\cdot}10^{-3}$ both give $\approx0$ for $R\leq0.32$. Against $R_0$ they disagree
strongly: frozen at $R=0.34,\ \lambda=3{\cdot}10^{-3}$ would give $f_{2+2}\approx0.3$--$0.4$,
the unfrozen network gives $0.000$. In the spec's terms: \textbf{representation adaptation
during $d4$ learning explains the difference} --- accessibility remains the governing variable,
but it must be evaluated on the final, self-consistently adapted representation.

\begin{figure}[h!]\centering
\includegraphics[width=0.58\textwidth]{phase_plane.pdf}\vspace{-4pt}
\caption{Phase-plane trajectories $R(t)$ vs.\ $f_{2+2}(t)$ for representative conditions
(seed 0; $\circ$ start, $\blacksquare$ end). Cold colors start from scratch ($R_0=0.19$) and
drift right; warm colors start pretrained ($R_0\approx0.33$) and, under decay, drift left ---
all end in the same sub-critical, low-$f_{2+2}$ band. $R(t)$ never crosses $R_c$.}
\label{fig:plane}
\end{figure}





\textbf{Conclusion.} In this task the compositional $2{+}2$ phase is reachable only with an
externally pinned (frozen or rescaled) representation. Left to its own dynamics the network
self-regulates to $R<R_c$ --- corner units cannot help emitting the singletons that dilute
their own pair advantage --- and weight decay, applied globally, reinforces the direct route by
shrinking the parent amplitudes themselves. A compositional phase under joint training would
require either parents whose production does not carry singleton by-products, deeper targets
(where the direct route's $\eta\propto Q_1/2^p$ collapses faster), or decay that spares the
representation-carrying layers.

\vspace{2pt}\footnotesize
Artifacts: \texttt{gdrive:parity/analysis\_output/accessibility\_unfrozen/}
(\texttt{results.csv}, \texttt{trajectories.csv}, \texttt{pretrain.json}, figures incl.\
\texttt{trajectories\_panels.pdf}, this report); code:
\texttt{analysis/accessibility/run\_unfrozen\_experiment.py}. Companions:
\texttt{accessibility\_report.pdf}, \texttt{accessibility\_wd\_report.pdf}.
\end{document}
